% EDITION: citation-based shortening, 4 October 2026.
% Earlier audits refer to the source edition; this revision has not been independently audited.
\section{P01: Well-posedness and bounded absorbing sets}
\label{sec:P01}

\subsection{Domain, hypotheses, and set terminology}
Use exactly the demand-modulated vector field in
\eqref{eq:canonical-dm}, denoted here by $F=(F_i)_{i\in\Zn}$.
The canonical assumptions imply the weaker hypotheses
\begin{equation}
\label{eq:P01-weak}
 n\ge1,\quad q\in\Z\pmod n,\quad
 a\in\R,\quad k>0,\quad c,g,h\ge0.
\tag{P01-W}
\end{equation}
We prove the main results under \eqref{eq:P01-weak}, without changing the
canonical contract for later modules. Coincident offsets retain every
term of the defining equation. Write
\[
 \Domain=(-1,\infty)^n,\qquad K_r=[0,r]^n\ (r\ge0),\qquad
 X(t)=\max_{i\in\Zn}x_i(t),\qquad s_+=\max\{s,0\}.
\]
For a continuous real function $v$, define
\[
 D^+v(t)=\limsup_{\tau\downarrow0}
             \frac{v(t+\tau)-v(t)}{\tau}.
\]
Distances use the maximum norm; for $x\in\Orth$,
\begin{equation}
\label{eq:P01-distance}
 \operatorname{dist}_\infty(x,K_r)=(\max_i x_i-r)_+.
\end{equation}
A set is \emph{positively invariant} if solutions starting there remain
there for all forward times for which they exist. It \emph{absorbs bounded
subsets of $\Orth$} if every bounded $E\subset\Orth$ has a finite $T_E$
such that all solutions starting in $E$ lie in the set for every
$t\ge T_E$. Attraction means convergence of the distance to zero;
uniform attraction of bounded sets means convergence uniform over each
such $E$. We will also establish stronger absorption bounds with a
single time valid for the entire orthant.

\subsection{Finite maxima, local solutions, and coordinate boundaries}
\label{slot:P01-1}

\begin{lemma}[Finite maxima, including ties]
\label{lem:P01-max}
For finitely many continuously differentiable functions $y_j$, set
$v(t)=\max_j y_j(t)$ and $I(t)=\{j:y_j(t)=v(t)\}$.
Then $v$ is locally Lipschitz, its right derivative exists at every
time with a right neighborhood in the interval, and
\begin{equation}
\label{eq:P01-max}
 D^+v(t)=\max_{j\in I(t)}\dot y_j(t).
\end{equation}
\end{lemma}
\begin{proof}
Bounded derivatives on compact time intervals make the finite maximum
Lipschitz. At a fixed $t$, every nonmaximizing index has a positive gap
from $v(t)$. Finiteness and continuity ensure that none of these indices
maximizes at $t+\tau$ for sufficiently small $\tau>0$. Uniformly over
$I(t)$,
$y_j(t+\tau)=v(t)+\tau\dot y_j(t)+o(\tau)$.
Taking the maximum and dividing by $\tau$ proves the assertion.
\end{proof}

\begin{proposition}[Local well-posedness and nonnegative invariance]
\label{prop:P01-local}
For any finite real coefficients, $F$ is $C^\infty$ on $\Domain$.
Every $x^0\in\Domain$ has a unique maximal classical forward solution
$x:[0,T_{\max})\to\Domain$, with continuous dependence on initial data.
If $c\ge0$ and $x^0\in\Orth$, then $x(t)\in\Orth$ for $t<T_{\max}$,
regardless of the signs of the other coefficients.
\end{proposition}
\begin{proof}
The denominators $1+x_i$ and $1+x_{i+q}$ are strictly positive on
$\Domain$. Each component of $F$ is a sum of polynomials and rational
functions with nonvanishing denominators, hence is smooth and locally
Lipschitz. Standard local ODE existence, uniqueness, and maximal
continuation theory on an open domain applies.

On each coordinate face,
\begin{equation}
\label{eq:P01-face}
 x_i=0,\quad x\in\Orth\quad\Longrightarrow\quad F_i(x)=c x_{i-1}\ge0.
\end{equation}
To treat tangent faces rigorously, fix $T<T_{\max}$ and let $x^+(t)$
be the coordinatewise positive part of $x(t)$. The curves $x,x^+$ and
their connecting line segments lie in a compact rectangle in $\Domain$.
Choose a Lipschitz constant $L$ for $F$ on this rectangle, in the
maximum norm, and put
\[
 V(t)=\max\{0,-x_0(t),\ldots,-x_{n-1}(t)\}
      =\norm{x(t)-x^+(t)}_\infty.
\]
If $V>0$ and $-x_i=V$, then $x_i^+=0$, so
$F_i(x^+)\ge0$ and $-F_i(x)\le |F_i(x)-F_i(x^+)|\le LV$.
If $V=0$, any active coordinate has $x_i=0$ and $-F_i(x)\le0$;
the constant zero is also active. Lemma~\ref{lem:P01-max} gives
$D^+V\le LV$. Since $V$ is Lipschitz, this inequality holds for its
ordinary derivative almost everywhere. Integration of
$(e^{-Lt}V)' \le0$, with $V(0)=0$, gives $V=0$ on $[0,T]$.
The choice of $T$ was arbitrary.
\end{proof}

\begin{proposition}[Persistence and propagation of positivity]
\label{prop:P01-positive}
For $c\ge0$ and nonnegative initial data, every initially positive
coordinate remains strictly positive on $[0,T_{\max})$. In particular,
strictly positive initial data stay strictly positive. If $c>0$ and
$x^0\ne0$, all coordinates are strictly positive at every
$0<t<T_{\max}$. If $c=0$, initially zero coordinates remain zero.
The zero initial state remains identically zero.
\end{proposition}
\begin{proof}
Along a nonnegative solution write
\[
 \dot x_i=b_i(t)x_i+d_i(t),\qquad
 b_i=a-kx_i-\frac{g}{1+x_{i+q}}-h x_{i-q},\qquad
 d_i=\frac{c x_{i-1}}{1+x_i}\ge0.
\]
These coefficients are continuous on every compact time interval in
the maximal interval. The integrating-factor formula gives
\begin{equation}
\label{eq:P01-positive}
 x_i(t)=x_i^0 e^{\int_0^t b_i(s)\,ds}
       +\int_0^t e^{\int_s^t b_i(r)\,dr}d_i(s)\,ds.
\end{equation}
This proves persistence of an initially positive coordinate.
If $c>0$, a predecessor positive at every $s>0$ supplies a strictly
positive integrand to its successor at every $s>0$, making that successor
positive at every $t>0$. At most $n-1$ iterations along the step-one
cycle reach all nodes, independently of $q$. For $c=0$ the integral
vanishes, proving preservation of zero coordinates. Finally $F(0)=0$
and uniqueness proves the zero-data assertion.
\end{proof}

\subsection{Maximum comparison and the exact comparison radius}
\label{slot:P01-2}
Under \eqref{eq:P01-weak}, define
\[
 b(s)=a-ks+\frac{c}{1+s},\qquad H(s)=s b(s)\quad(s\ge0),
\]
and set
\begin{equation}
\label{eq:P01-radius}
 X_{\max}=R:=\max\left\{0,\frac{a-k+\sqrt{(a+k)^2+4kc}}{2k}\right\}.
\end{equation}
This radius is defined by a scalar upper comparison; it is not asserted
to be the smallest invariant cube of the full vector field.

\begin{lemma}[Dini inequality and scalar root]
\label{lem:P01-bound}
Every nonnegative solution under \eqref{eq:P01-weak} satisfies
\begin{equation}
\label{eq:P01-dini}
 D^+X=\max_{i:x_i=X}F_i(x)
       \le H(X)-\frac{gX}{1+X}\le H(X).
\end{equation}
If $a+c>0$, then $R>0$ is the unique positive zero of $b$, and
\begin{equation}
\label{eq:P01-polynomial}
 kR^2+(k-a)R-(a+c)=0.
\end{equation}
If $a+c\le0$, then $R=0$ and $H(s)\le-ks^2$ for $s\ge0$.
For $R>0$ and $s\ge R$,
\begin{equation}
\label{eq:P01-factor}
 H(s)=s\left[-k(s-R)
    +c\left(\frac1{1+s}-\frac1{1+R}\right)\right]\le-ks(s-R).
\end{equation}
\end{lemma}
\begin{proof}
Lemma~\ref{lem:P01-max} includes every tied maximizing coordinate.
At $x_i=X$, nonnegative invariance gives
$0\le x_{i-1},x_{i+q},x_{i-q}\le X$. Consequently
\[
 \frac{c x_{i-1}}{1+X}\le\frac{cX}{1+X},\qquad
 -\frac{gX}{1+x_{i+q}}\le-\frac{gX}{1+X},\qquad
 -hXx_{i-q}\le0.
\]
The middle inequality uses the negative prefactor and $g\ge0$.
Adding proves \eqref{eq:P01-dini}; no independence of neighbor
coordinates is required.

Now $b'(s)=-k-c/(1+s)^2<0$, $b(0)=a+c$, and $b(s)\to-\infty$.
If $a+c>0$, multiply $b(R)=0$ by $1+R$ to obtain
\eqref{eq:P01-polynomial}. Its constant term is negative and its
discriminant is $(a+k)^2+4kc$, so its positive root is
\eqref{eq:P01-radius}. If $a+c\le0$, then $k-a>0$ and
\[
 (a+k)^2+4kc-(k-a)^2=4k(a+c)\le0.
\]
Thus the raw root in \eqref{eq:P01-radius} is nonpositive, $R=0$,
and $b(s)\le a+c-ks\le-ks$. Subtracting $b(R)=0$ when $R>0$
gives \eqref{eq:P01-factor}; its inequality follows from $c\ge0$
and $s\ge R$.
\end{proof}

\begin{theorem}[Global continuation and invariant attracting cubes]
\label{thm:P01-global}
Under \eqref{eq:P01-weak}, every $x^0\in\Orth$ has a unique classical
solution for all $t\ge0$. The scalar solution
\begin{equation}
\label{eq:P01-scalar}
 \dot u=H(u),\qquad u(0)=X(0)
\end{equation}
is global and nonnegative, and
\begin{equation}
\label{eq:P01-globalbound}
 0\le X(t)\le u(t)\le\max\{X(0),R\}.
\end{equation}
Every $K_r$ with $r\ge R$ is positively invariant. Furthermore,
\begin{align}
 \operatorname{dist}_\infty(x(t),K_R)
   &\le e^{-kRt}(X(0)-R)_+ &&(R>0), \label{eq:P01-distance-rate}\\
 X(t)&\le\frac{X(0)}{1+kX(0)t} &&(R=0). \label{eq:P01-zero-rate}
\end{align}
There are also initial-state-independent bounds, for $t>0$:
\begin{align}
 X(t)&\le\frac{R}{1-e^{-kRt}} &&(R>0), \label{eq:P01-comingdown}\\
 X(t)&\le\frac1{kt} &&(R=0). \label{eq:P01-comingdown-zero}
\end{align}
Thus $K_R$ attracts all nonnegative solutions uniformly on bounded
initial sets, and indeed uniformly over the entire orthant as
$t\to\infty$. Every larger cube $K_r$, $r>R$, absorbs the entire
orthant by the common time
\begin{equation}
\label{eq:P01-outer-time}
 T_r=\begin{cases}
 \displaystyle\frac1{kR}\log\frac{r}{r-R},&R>0,\\[4pt]
 \displaystyle\frac1{kr},&R=0.
 \end{cases}
\end{equation}
\end{theorem}
\begin{proof}
The scalar field $H$ is smooth on $(-1,\infty)$, $H(0)=0$, and
$H(s)\le0$ for $s\ge R$. Its solution starting nonnegative stays in
$[0,M]$, where $M=\max\{u(0),R\}$: uniqueness prevents crossing zero,
and the upper endpoint has nonpositive field. The latter barrier also
follows by the positive-part argument below applied to $(u-M)_+$,
using $H(M)\le0$. Compact containment gives global scalar continuation.

We justify the vector comparison without assuming differentiability of
$X$ everywhere. On a compact subinterval of the vector maximal interval,
let $L$ be a Lipschitz constant for $H$ on an interval containing $X$
and $u$. The Lipschitz function
\[
 W=(X-u)_+=\max\{0,x_i-u:i\in\Zn\}
\]
satisfies $D^+W\le H(X)-H(u)\le LW$ where $W>0$.
At $W=0$, any active coordinate has $x_i=u=X$, unless the constant
zero is the only active function; Lemma~\ref{lem:P01-max} and
\eqref{eq:P01-dini} give $D^+W\le0$. Integrating the almost-everywhere
inequality, as in Proposition~\ref{prop:P01-local}, yields $W=0$.
This proves \eqref{eq:P01-globalbound} on the maximal interval.

If $T_{\max}<\infty$, the vector trajectory stays in $K_M$, a compact
subset of $\Domain$ at distance at least one from every singular
hyperplane. The field is bounded on a compact neighborhood in $\Domain$,
hence $\dot x$ is bounded. Thus $x(t)$ has a limit in $K_M$ as
$t\uparrow T_{\max}$. Local existence at that limit and uniqueness
extend the solution past $T_{\max}$, a contradiction. This proves
global continuation and absence of finite-time blow-up. If
$X(0)\le r$ and $r\ge R$, \eqref{eq:P01-globalbound} proves invariance.

For $R>0$, \eqref{eq:P01-factor} gives
$\dot u\le-ku(u-R)\le-kR(u-R)$ where $u\ge R$. It follows that
$(u(t)-R)_+\le e^{-kRt}(u(0)-R)_+$, including solutions initially at
or below $R$, and hence \eqref{eq:P01-distance-rate}. If $u(0)>R$,
set $w=1-R/u$ while $u>R$. Then
\[
 \dot w=\frac{R\dot u}{u^2}\le-kR w,\qquad 0<w(0)<1.
\]
Solving this inequality gives \eqref{eq:P01-comingdown} for $u$ and
therefore for $X$. If $u$ reaches $R$, invariance supplies the same
bound afterwards; if $u(0)\le R$, it is immediate.
For $R=0$, integrating $(1/u)'\ge k$ where $u>0$, using
$\dot u\le-ku^2$, proves \eqref{eq:P01-zero-rate} and
\eqref{eq:P01-comingdown-zero}; zero initial data follow by uniqueness.

Attraction follows from these estimates and \eqref{eq:P01-distance}.
For $R>0$ the distance is at most $R/(e^{kRt}-1)$ independently of
initial data; for $R=0$ use $1/(kt)$. Substitution of $T_r$ into
\eqref{eq:P01-comingdown} or \eqref{eq:P01-comingdown-zero} gives
$X(T_r)\le r$, and invariance supplies all subsequent times.
\end{proof}

\subsection{Finite-time absorption at the comparison radius}
\label{slot:P01-3}

\begin{theorem}[Exact absorption criterion for $K_R$]
\label{thm:P01-absorption}
Assume \eqref{eq:P01-weak}. If $R=0$, no initially nonzero nonnegative
solution enters $K_R=\{0\}$ at finite time. This cube is invariant and
globally attracting but not absorbing.

If $R>0$, the following are equivalent:
\begin{enumerate}[label=(\roman*),leftmargin=*]
\item $K_R$ absorbs every nonnegative solution in finite time.
\item $K_R$ absorbs bounded subsets of $\Orth$.
\item One finite absorption time is valid for the entire orthant.
\item The coefficients and offset satisfy
\begin{equation}
\label{eq:P01-absorption-criterion}
 g>0\quad\hbox{or}\quad
 [\,h>0\ \hbox{and}\ (c>0\ \hbox{or}\ q\equiv0\pmod n)\,].
\end{equation}
\end{enumerate}
Canonically $R>0$, $g>0$, and the entire orthant enters $K_R$ by
\begin{equation}
\label{eq:P01-canonical-time}
 T_*=\frac{\log2}{kR}+\frac{1+R}{g}.
\end{equation}
If $g>0$ and $X(0)\le M$, an alternative entry-time bound is
\begin{equation}
\label{eq:P01-bounded-time}
 T_M=\frac{(M-R)_+(1+R)}{gR}.
\end{equation}
\end{theorem}
\begin{proof}
When $R=0$, Proposition~\ref{prop:P01-positive} ensures that some
coordinate of any initially nonzero solution stays positive at each
finite time.

Now suppose $R>0$ and set $t_0=(\log2)/(kR)$. By
\eqref{eq:P01-comingdown}, every trajectory is in the invariant cube
$K_{2R}$ by $t_0$. If $g>0$, then whenever $X\ge R$,
\begin{equation}
\label{eq:P01-strict-gap}
 D^+X\le-\frac{gX}{1+X}\le-\frac{gR}{1+R}=:-\delta<0,
\end{equation}
using $H(X)\le0$. Absolute continuity permits integration before
the first entry into $X\le R$. If $X(0)\le M$, a longer time outside
than $(M-R)_+/\delta$ would force $X\le R$, a contradiction.
Invariance gives all later times. Applying this after $t_0$ with
$M=2R$ proves \eqref{eq:P01-canonical-time}.

If $g=0$, $h>0$, and $q\equiv0\pmod n$, a maximizing coordinate
has $x_{i-q}=X$. Thus $D^+X\le H(X)-hX^2\le-hR^2$ when $X\ge R$,
giving common absorption by $t_0+1/(hR)$.

For the remaining sufficient case $g=0$, $h>0$, $c>0$, we give a
coordinatewise quantitative argument. Put $M=2R$ and fix any
$\tau>0$. On $K_M$,
\begin{equation}
\label{eq:P01-propagation-bound}
 \dot x_i\ge-\beta x_i+\eta x_{i-1},\qquad
 \beta=(-a)_++(k+h)M+g,\quad \eta=\frac{c}{1+M}>0.
\end{equation}
Indeed $a\ge-(-a)_+$, $kx_i\le kM$,
$g/(1+x_{i+q})\le g$, $h x_{i-q}\le hM$, and
$c/(1+x_i)\ge c/(1+M)$. On $[s,s+\tau]$ in this cube, the transformed
variables $z_i(r)=e^{\beta(r-s)}x_i(r)$ satisfy
\[
 z_i(r)\ge x_i(s)+\eta\int_s^r z_{i-1}(v)\,dv.
\]
Repeated substitution, keeping the remaining nested integral
nonnegative, yields
\begin{equation}
\label{eq:P01-propagation-series}
 x_i(s+\tau)\ge e^{-\beta\tau}
       \sum_{j=0}^{n-1}\frac{(\eta\tau)^j}{j!}x_{i-j}(s).
\end{equation}
A $j$-fold integral of a constant over the time simplex equals
$\tau^j/j!$, which gives the displayed coefficients.

For a trajectory that has not entered $K_R$ by time $t\ge t_0+\tau$,
we have $X(t-\tau)>R$. The indices $i-j$, $0\le j\le n-1$, cover
the cycle once. Thus \eqref{eq:P01-propagation-series} gives, for every $i$,
\begin{equation}
\label{eq:P01-positive-floor}
 x_i(t)\ge d:=R e^{-\beta\tau}
        \min_{0\le j\le n-1}\frac{(\eta\tau)^j}{j!}>0.
\end{equation}
At such times, every maximizing coordinate satisfies
$F_i(x)\le H(X)-hXx_{i-q}\le-hRd$.
Starting from $X\le2R$ at $t_0+\tau$, integration proves entry by
$t_0+\tau+1/(hd)$, or earlier. We have proved (iv)$\Rightarrow$(iii).

The implications (iii)$\Rightarrow$(ii)$\Rightarrow$(i) are immediate.
If (iv) fails and $g=h=0$, choose all initial coordinates equal to
$s_0>R$. Direct substitution verifies the solution $x_i(t)=s(t)$
with $\dot s=H(s)$. It stays above $R$ at each finite time:
a first meeting with the constant scalar solution $R$ would violate
uniqueness, also backwards near that meeting. Hence it never enters
$K_R$. The other failure case is $g=0$, $h>0$, $c=0$,
$q\not\equiv0\pmod n$. Here $R=a/k>0$. Choose one initial coordinate
$x_j(0)>R$ and all others zero. The zero coordinates stay zero and
$j-q\ne j$, so $\dot x_j=a x_j-kx_j^2$. It likewise stays strictly
above $a/k$ at every finite time. These cases exhaust failure of (iv)
and contradict (i), completing the equivalence.
\end{proof}

\subsection{Canonical conclusions and weaker-assumption limits}
\label{slot:P01-4}

\begin{corollary}[Canonical foundation]
\label{cor:P01-canonical}
For canonical $n\ge5$, $1\le q\le n-1$ and $a,k,c,g,h>0$, every
nonnegative initial state has a unique global nonnegative solution.
Every initially positive coordinate stays positive, and every nonzero
initial state becomes strictly positive in every coordinate at each
positive time. The cube $[0,X_{\max}]^n$, with
\[
 X_{\max}=\frac{a-k+\sqrt{(a+k)^2+4kc}}{2k}>0,
\]
is positively invariant, globally attracting, and finite-time absorbing
for the entire orthant with the common bound
\eqref{eq:P01-canonical-time}. All larger cubes are also invariant
and absorbing. No condition on $\gcd(n,q)$ is required.
\end{corollary}
\begin{proof}
The canonical assumptions imply \eqref{eq:P01-weak}, $a+c>0$,
$c>0$ and $g>0$. Apply the preceding results.
\end{proof}

\begin{remark}[Archived radius and the meaning of stability]
The blueprint's candidate radius
\[
 B=\frac{a+\sqrt{a^2+4kc}}{2k}
\]
is the larger real root of $as-ks^2+c=0$, and is nonnegative. If $c>0$, then $B>0$ and
$H(B)=-c/(1+B)<0$, so $B>R$. Its cube is therefore also invariant
and absorbing by Theorem~\ref{thm:P01-global}, even when $g=h=0$.
The smaller radius $R$ retains $cX/(1+X)$ instead of replacing it
by $c$. If $c=0$, then $B=R$ and the exact-radius criterion applies.
Neither radius is claimed optimal after all $g,h,q$ information is used.

Invariance concerns remaining inside; absorption concerns finite-time
entry; attraction alone permits approaching from outside forever.
The distance estimates also show Lyapunov stability of $K_R$ as a set,
relative to $\Orth$, because distance never exceeds its initial value.
Under the convention that asymptotic stability of a set means Lyapunov
stability plus attraction, $K_R$ is globally asymptotically stable as
a set. This is not convergence to a particular point, nor a claim that
the cube is a minimal global attractor or is invariant in both time
directions.
\end{remark}

\begin{proposition}[Global existence when $k=0$ is allowed]
\label{prop:P01-kzero}
Global nonnegative existence and the positivity conclusions still hold
for $k\ge0$, $a\in\R$, $c,g,h\ge0$, any $n\ge1$, and any residue $q$.
For $k=0$, no bounded absorbing cube is guaranteed over this parameter
class.
\end{proposition}
\begin{proof}
Local existence and nonnegative invariance did not require $k>0$.
Discarding nonpositive losses at a maximizing coordinate gives
\[
 D^+X\le aX+\frac{cX}{1+X}\le(a_++c)X.
\]
The Lipschitz/Dini argument yields $X(t)\le X(0)e^{(a_++c)t}$.
If the maximal time were finite, this would confine the trajectory
to a compact subset of $\Domain$, and the earlier continuation
argument would extend it. Positivity follows from
Proposition~\ref{prop:P01-positive}. For the last assertion take
$a>0$, $k=c=g=h=0$: an initially positive coordinate equals
$x_i(0)e^{at}$ and is unbounded.
\end{proof}
